\documentclass[11pt]{amsart}

\usepackage[T1]{fontenc}
\usepackage[utf8]{inputenc}
\usepackage{lmodern}
\usepackage{microtype}
\setlength{\emergencystretch}{2em}
\usepackage[margin=1.15in]{geometry}
\usepackage{amsmath,amssymb,mathtools}
\usepackage{mathrsfs}
\usepackage{amsthm}
\usepackage{booktabs}
\usepackage{enumitem}
\usepackage{xcolor}
\usepackage[colorlinks=true,linkcolor=blue!55!black,urlcolor=blue!55!black]{hyperref}

\newtheorem{theorem}{Theorem}[section]
\newtheorem{proposition}[theorem]{Proposition}
\newtheorem{lemma}[theorem]{Lemma}
\newtheorem{corollary}[theorem]{Corollary}
\newtheorem{definition}[theorem]{Definition}
\newtheorem{example}[theorem]{Example}
\newtheorem{remark}[theorem]{Remark}

\newcommand{\C}{\mathbb C}
\newcommand{\End}{\operatorname{End}}
\newcommand{\Hom}{\operatorname{Hom}}
\newcommand{\rank}{\operatorname{rank}}
\newcommand{\id}{\operatorname{id}}

\title[Resonance-Marked Naturality]{Resonance-Marked Naturality and Homotopy-Tilt Non-Invariance\\
in the Inverse-Leibniz Cubic Case}
\author{Akari H.}
\date{Paper II draft v0.01 --- September 2026}

\begin{document}

\begin{abstract}
We isolate the chain-level datum that remains after the fixed cubic inverse-Leibniz deformation problem is compared across homotopy splittings.  A normalized recurrent envelope carries an affine action of admissible homotopy tilts.  In the cubic model the seven-dimensional recurrent operator has a one-dimensional normalized direction generated by \(E\), with \(A(E)=-2E\), while its complementary spectral sector can be changed arbitrarily by internal tilts.  Thus the quintic factor
\[
q(S)=S^5-4S^4+12S^3-32S^2+80S-192
\]
and the fixed-contraction polynomial \((S+2)q(S)\) are not unmarked deformation invariants.  The surviving object is instead the primitive-free resonant boundary-response pair
\[
(\mathcal L_{\mathrm{res}},\rho_{\mathrm{res}}),
\qquad
\mathcal L_{\mathrm{res}}=\C(\xi\smile\xi),
\]
which is naturally transported by strict comparisons preserving the resonance marking and its primitive quotient.  We prove this positive naturality statement and give a small strict DGLA quasi-isomorphism that cancels the contractible resonant pair while preserving the underlying formal deformation behavior.  Consequently the pair is a resonance-marked chain-level invariant, but not an ordinary derived invariant.  This is the stopping point of Paper II; the general triangular theory with arbitrary spectral multiplicities is deferred.
\end{abstract}

\maketitle
\tableofcontents

\section{Scope and main conclusions}

The inverse-Leibniz cubic calculation produces two kinds of output.  The first kind belongs to a chosen contraction: a controller box, a recurrent matrix, a cyclic source generator, and a polynomial such as
\[
m_A(S)=(S+2)q(S).
\]
The second kind is forced by the original cochain identities together with a normalized boundary direction.  The purpose of this paper is to separate these two layers and to identify the correct comparison statement.

The main conclusions are:

\begin{enumerate}[label=(\roman*)]
\item the exact recurrent operator is not invariant under admissible homotopy tilts;
\item the normalized \(E\)-rail and its \(-2\) response are invariant inside the fixed envelope;
\item the primitive-free resonant boundary-response pair is natural under strict resonance-faithful comparisons;
\item ordinary DGLA quasi-isomorphism does not preserve that pair, even when the formal Maurer--Cartan deformation behavior is unchanged.
\end{enumerate}

The paper therefore ends with the boxed scope statement
\[
\boxed{\text{strict/resonance-faithful naturality}\;+
       \text{ordinary quasi-isomorphism no-go}.}
\]

No claim is made here for arbitrary triangular systems with spectral multiplicity vector
\[
m_A(S)=\prod_\lambda(S-\lambda)^{e_\lambda}.
\]

\section{The fixed cubic recurrent envelope}

Let
\[
C^1\xrightarrow{d}C^2\xrightarrow{d}C^3
\]
be the relevant constrained Hochschild complex, and let
\[
D_\alpha=[\alpha,-]:C^1\longrightarrow C^2
\]
be the operator determined by the reduced tangent direction \(\alpha\).  A selected partial contraction supplies a finite-dimensional recurrent envelope
\[
\mathcal R\subset C^1,
\qquad
A=-hD_\alpha\big|_{\mathcal R}:\mathcal R\longrightarrow\mathcal R.
\]

The exact Part II computation gives
\[
\dim\mathcal R=7,
\qquad
\rank(D_\alpha|_{\mathcal R})=7,
\]
and
\[
D_\alpha(\mathcal R)\cap Z^2
=D_\alpha(\mathcal R)\cap B^2
=\C D_\alpha E.
\]
The normalized direction satisfies
\[
D_\alpha E=2dE=-2\,\xi\smile\xi,
\qquad
A(E)=-2E.
\]

For the fixed contraction, the characteristic polynomial is
\[
\chi_A(S)=(S+2)^2q(S),
\qquad
q(S)=S^5-4S^4+12S^3-32S^2+80S-192.
\]
This formula is exact, but its status is deliberately local: it describes the selected partial contraction before quotienting by homotopy tilts.

\section{The exact homotopy-tilt action}

\begin{definition}[Normalized source directions]
Put
\[
F=D_\alpha(\mathcal R),
\qquad
F_{\mathrm{norm}}=F\cap Z^2,
\qquad
\overline F=F/F_{\mathrm{norm}}.
\]
The normalized internal tilt group is
\[
\mathbb T_{\mathcal R}=\Hom(\overline F,\mathcal R).
\]
\end{definition}

For \(\tau\in\mathbb T_{\mathcal R}\), write \(q_F:F\to\overline F\) for the quotient map and define
\[
h_\tau(q)=h(q)+\tau(q_F(q)),
\qquad q\in F.
\]
The normalized directions are fixed because \(q_F\) vanishes on \(F_{\mathrm{norm}}\).  The residual source changes by
\[
r_\tau(q)=q-dh_\tau(q),
\qquad
dr_\tau(q)=dq.
\]
Thus the degree-three differential data of the nonclosed source quotient are unchanged.

\begin{proposition}[Exact operator action]
The recurrent operator changes according to
\[
A_\tau=A-\tau q_FD_\alpha|_{\mathcal R}.
\]
The map
\[
\mathbb T_{\mathcal R}
\longrightarrow
\{T\in\End(\mathcal R):T(E)=0\},
\qquad
\tau\longmapsto \tau q_FD_\alpha|_{\mathcal R},
\]
is an isomorphism in the cubic model.
\end{proposition}

\begin{proof}
The first identity is immediate from \(A=-hD_\alpha\) and the definition of \(h_\tau\).  Since \(D_\alpha|_{\mathcal R}\) has rank seven while \(F_{\mathrm{norm}}=\C D_\alpha E\), the quotient map \(q_FD_\alpha\) has kernel \(\C E\) and identifies
\[
\mathcal R/\C E\cong\overline F.
\]
Every endomorphism vanishing on \(E\) therefore factors uniquely through \(q_FD_\alpha\).
\end{proof}

\begin{corollary}[Cubic tilt orbit]
In the fixed seven-dimensional envelope,
\[
\dim\mathbb T_{\mathcal R}=6\cdot7=42,
\]
and
\[
\boxed{
\mathcal O_{\mathrm{tilt}}(A)
=\{B\in\End(\mathcal R):B(E)=-2E\}.}
\]
\end{corollary}

\begin{proof}
The dimension follows from \(\dim\overline F=6\) and \(\dim\mathcal R=7\).  The preceding proposition says precisely that all variations vanishing on \(E\) occur, while the normalized relation \(A(E)=-2E\) is fixed.
\end{proof}

\section{Spectral non-invariance and the surviving rail}

The tilt orbit contains the two exact representatives
\[
B_0=\operatorname{diag}(-2,0,0,0,0,0,0),
\qquad
B_3=\operatorname{diag}(-2,3,3,3,3,3,3).
\]
Their characteristic polynomials are
\[
\chi_{B_0}(S)=(S+2)S^6,
\qquad
\chi_{B_3}(S)=(S+2)(S-3)^6.
\]
Both operators agree with \(A\) on \(\C E\) and hence are exact normalized internal tilts.  We obtain:

\begin{theorem}[Homotopy-tilt spectral no-go]
The factor \(q(S)\), the second \(-2\) eigenline, the six-dimensional cyclic module generated by the fixed \(g_3(0)\), and the fixed-contraction recurrence \((S+2)q(S)\) are not invariants of the original cubic deformation problem under normalized internal homotopy tilts.  The greatest common divisor of the characteristic polynomials over the full internal tilt orbit is
\[
\gcd_{\tau}\chi_{A_\tau}(S)=S+2,
\]
and the same statement holds for minimal polynomials.
\end{theorem}

\begin{proof}
The two displayed representatives already change the complementary spectrum.  Conversely every operator in the orbit has the eigenrelation \(B(E)=-2E\), so \(S+2\) divides every characteristic and minimal polynomial.  Since the induced operator on \(\mathcal R/\C E\) is arbitrary, no nonconstant factor on that quotient can be common to the whole orbit.
\end{proof}

The surviving temporal response is therefore the normalized \(E\)-rail:
\[
\boxed{
D_\alpha A^nE=2(-2)^n dE
\qquad(n\geq0).}
\]
This is a chain-level boundary response, not the full spectrum of a chosen recurrent matrix.

\begin{remark}
The dimension seven in this section is an envelope dimension for the present internal tilt calculation.  It is not asserted to be invariant under an arbitrary change of partial contraction that changes the recurrent envelope itself.
\end{remark}

\section{Why the old \(g_3\) source is not canonical}

The fixed low-rail construction has
\[
g_3(0)=2\mathsf k-\mathsf p.
\]
At the first nontrivial source stage, the admissible triangular splitting tilt
\[
\mathsf p_t=\mathsf p+t u,
\qquad
C_{13,t}=C_{13}-tdu,
\]
leaves the original source equation and its differential-rank placement unchanged.  Taking \(u=E\) gives
\[
g_{3,t}(0)=g_3(0)-tE,
\]
and hence
\[
D_\alpha g_{3,t}(0)
=D_\alpha g_3(0)-tD_\alpha E.
\]
Since \(D_\alpha E\neq0\), the exact \(g_3\) source sequence is not determined by the original equations independently of the contraction.  Taking \(u=\mathsf k\) gives the stronger variation
\[
D_\alpha g_{3,t}(0)
=D_\alpha g_3(0)-t[\alpha,\mathsf k],
\qquad
d[\alpha,\mathsf k]\neq0,
\]
so this ambiguity is not merely a boundary ambiguity.

In contrast, the \(E\)-rail is forced by the original identities
\[
dE=-\xi\smile\xi,
\qquad
[\alpha,E]=-2\xi\smile\xi=2dE,
\]
and the normalized contraction condition \(h(dE)=E\).  Therefore
\[
A(E)=-2E
\]
is stable under the normalized internal action.

\section{The primitive-free resonant response object}

The \(E\)-dependent description is useful for computation but not optimal for comparison.  We now remove the primitive choice.

\begin{definition}[Resonant boundary-response pair]
Let
\[
\beta=\xi\smile\xi,
\qquad
\mathcal L_{\mathrm{res}}=\C\beta\subset B^2,
\]
and define
\[
\mathcal P_{\mathrm{res}}
=d^{-1}(\mathcal L_{\mathrm{res}})/Z^1.
\]
The differential induces an isomorphism
\[
\bar d:\mathcal P_{\mathrm{res}}\xrightarrow{\sim}\mathcal L_{\mathrm{res}}.
\]
\end{definition}

Let \(T_{\mathrm{red}}\) be the reduced tangent space.  For \(a\in T_{\mathrm{red}}\), assume
\[
[a,Z^1]=0,
\qquad
[a,d^{-1}(\mathcal L_{\mathrm{res}})]
\subseteq\mathcal L_{\mathrm{res}}.
\]
Then \(D_a=[a,-]\) descends to
\[
\bar D_a:\mathcal P_{\mathrm{res}}\longrightarrow\mathcal L_{\mathrm{res}}.
\]

\begin{definition}[Response operator and covector]
Define
\[
\Theta(a)=\bar D_a\bar d^{-1}
\in\End(\mathcal L_{\mathrm{res}}).
\]
Since \(\mathcal L_{\mathrm{res}}\) is one-dimensional, there is a unique scalar \(\rho_{\mathrm{res}}(a)\) such that
\[
\Theta(a)=\rho_{\mathrm{res}}(a)\id.
\]
Thus
\[
\rho_{\mathrm{res}}\in T_{\mathrm{red}}^\vee.
\]
In the normalized cubic coordinate,
\[
\rho_{\mathrm{res}}(\alpha)=2.
\]
\end{definition}

\section{Strict resonance-faithful naturality}

\begin{definition}[Resonance-faithful comparison]
Let two constrained DGLA models carry marked data
\[
(T_{\mathrm{red}},\mathcal L_{\mathrm{res}},\mathcal P_{\mathrm{res}},\bar d,\bar D)
\quad\text{and}\quad
(T'_{\mathrm{red}},\mathcal L'_{\mathrm{res}},\mathcal P'_{\mathrm{res}},\bar d',\bar D').
\]
A comparison consists of isomorphisms
\[
\Phi_T:T_{\mathrm{red}}\xrightarrow{\sim}T'_{\mathrm{red}},
\quad
\Phi_L:\mathcal L_{\mathrm{res}}\xrightarrow{\sim}\mathcal L'_{\mathrm{res}},
\quad
\Phi_P:\mathcal P_{\mathrm{res}}\xrightarrow{\sim}\mathcal P'_{\mathrm{res}},
\]
such that
\[
\bar d'\Phi_P=\Phi_L\bar d,
\qquad
\bar D'_{\Phi_Ta}\Phi_P=\Phi_L\bar D_a.
\]
If the comparison comes from a strict DGLA isomorphism, these identities follow from preservation of the differential and bracket together with preservation of the resonance marking.
\end{definition}

\begin{theorem}[Naturality of the marked response]
Under a resonance-faithful comparison,
\[
\boxed{
\Theta'(\Phi_Ta)=\Phi_L\Theta(a)\Phi_L^{-1}.}
\]
Consequently,
\[
\boxed{
\rho'_{\mathrm{res}}(\Phi_Ta)=\rho_{\mathrm{res}}(a),}
\]
or equivalently
\[
\rho'_{\mathrm{res}}=(\Phi_T^{-1})^\vee\rho_{\mathrm{res}}.
\]
\end{theorem}

\begin{proof}
Using the two commuting identities gives
\[
\begin{aligned}
\Theta'(\Phi_Ta)
&=\bar D'_{\Phi_Ta}(\bar d')^{-1}\\
&=\Phi_L\bar D_a\Phi_P^{-1}\Phi_L^{-1}\\
&=\Phi_L\Theta(a)\Phi_L^{-1}.
\end{aligned}
\]
The resonant line is one-dimensional, so conjugation leaves its scalar unchanged.  The covector formula is the same statement in dual form.
\end{proof}

\begin{remark}
Rescaling the normal generator changes \(\beta\) and hence the chosen basis of \(\mathcal L_{\mathrm{res}}\), but the conjugation formula cancels that scalar change.  The line and the response operator, rather than a fixed boundary vector, are the intrinsic marked objects.
\end{remark}

\section{Ordinary quasi-isomorphism no-go}

\begin{theorem}[Contractible-pair cancellation]
There are strict DGLA quasi-isomorphisms that preserve the formal Maurer--Cartan deformation behavior but erase the marked resonant boundary-response pair.
\end{theorem}

\begin{proof}
Let \(\mathfrak g\) have degree-one generators \(\alpha,\xi,E\) and degree-two generator \(\beta\), with
\[
d\alpha=d\xi=0,
\qquad
dE=\beta,
\]
and the only nonzero brackets
\[
[\xi,\xi]=2\beta,
\qquad
[\alpha,E]=2\beta.
\]
Let \(\mathfrak h\) be the abelian DGLA on degree-one generators \(\alpha,\xi\), and let
\[
p:\mathfrak g\longrightarrow\mathfrak h
\]
kill \(E\) and \(\beta\) while fixing \(\alpha\) and \(\xi\).  The pair \(E\to\beta\) is contractible, so \(p\) is a strict quasi-isomorphism.  Both models have two-dimensional degree-one cohomology and vanishing degree-two cohomology.

For \(\gamma=u\alpha+v\xi+wE\), the Maurer--Cartan equation in \(\mathfrak g\) is
\[
w+v^2+2uw=0,
\]
so formally \(w=-v^2/(1+2u)\) and the deformation space is the same smooth two-parameter germ as for \(\mathfrak h\).

Nevertheless, in \(\mathfrak g\) one has
\[
\mathcal L_{\mathrm{res}}=\C\beta,
\qquad
\rho_{\mathrm{res}}(\alpha)=2,
\]
whereas in \(\mathfrak h\) the product \([\xi,\xi]\) vanishes and
\[
\mathcal L'_{\mathrm{res}}=0.
\]
Thus the ordinary quasi-isomorphism erases precisely the contractible chain-level datum that the marked theory is designed to retain.
\end{proof}

\begin{corollary}
The pair \((\mathcal L_{\mathrm{res}},\rho_{\mathrm{res}})\) is not an invariant of the ordinary DGLA derived category or of the associated formal moduli problem alone.
\end{corollary}

\section{Verification boundary and paper closeout}

The finite cubic statements in this paper are supported by exact rational linear algebra.  The current Part II verification family establishes:

\begin{itemize}
\item the support/incidence model and the \(35\)-dimensional \(U_{16}\) controller image;
\item the \(19\)-dimensional forced domain and \(7\)-dimensional recurrent image;
\item the comparison behavior of rigid germs;
\item the \(42\)-dimensional internal tilt module and the exact orbit formula;
\item the two contrasting characteristic polynomials proving spectral non-invariance.
\end{itemize}

The v0.53 source-behavior record gives the explicit \(E\)- and \(g_3\)-rail identities used above.  Its current standalone verifier still imports an absent historical helper, so the reproducibility package must bundle that helper or replace it with a self-contained check before submission.  This packaging item does not enlarge the mathematical scope of the paper.

The following questions are intentionally excluded from Paper II:

\begin{enumerate}[label=(\arabic*)]
\item a general theorem for arbitrary \(m_A(S)=\prod_\lambda(S-\lambda)^{e_\lambda}\);
\item a proof that all admissible partial contractions have isomorphic rigid germs;
\item contraction-independent classification of the full \(U_{16}\) envelope;
\item a canonical gauge-fixing that recovers the richer \(q\)-spectrum;
\item simultaneous deformation of the algebra \(A\) and the tensor \(B\).
\end{enumerate}

The final statement of Paper II is therefore:
\[
\boxed{
(\mathcal L_{\mathrm{res}},\rho_{\mathrm{res}})
\text{ is natural for resonance-faithful marked comparisons,}
}
\]
but
\[
\boxed{
(\mathcal L_{\mathrm{res}},\rho_{\mathrm{res}})
\text{ is not ordinary quasi-isomorphism invariant.}}
\]

\end{document}
